\emph{Setting.} Gravitational fields $(\hb, \Omega)$ on flat spacetime are as in Definition \claimref{obj.synthesis.gtg-flat-spacetime-gauge-fields}; only clauses (i) to (e) are used (the fields, the covariant derivative, the field strength, the derived metric, the gauge freedoms), never the field equation (f). At first order about $\hb = \id$, $\Omega = 0$ write $\hb^{-1}(a) = a + \tfrac12 h(a)$ with $h$ symmetric (the antisymmetric part of a first-order $\hb$ is a rotation gauge direction), so $g(a, b) = a \cdot b + h(a) \cdot b$; the torsion-free connection and its curvature are read through Fact \claimref{synthesis.imported-gtg-dictionary}. For the polynomial field $h(x) = \tfrac12 (p \cdot x)^{2}h_{0}$ the second derivatives of $g$ are $(a \cdot p)(b \cdot p)\,h_{0}$ and the first-order curvature is constant. In components the first-order Riemann tensor of $g_{ab} + h_{ab}$ is $\tfrac12(\partial_{b}\partial_{c}h_{ad} + \partial_{a}\partial_{d}h_{bc} - \partial_{a}\partial_{c}h_{bd} - \partial_{b}\partial_{d}h_{ac})$ (a labelled index reading), and pairing it with $c \wedge d$ gives, for $B = a \wedge b$, the bivector $\tfrac12[(a \cdot p)\,p \wedge h_{0}(b) - (b \cdot p)\,p \wedge h_{0}(a)] = -\tfrac12\,p \wedge h_{0}(B \cdot p)$, since $(a \wedge b) \cdot p = a(b \cdot p) - b(a \cdot p)$ and $(c \wedge d) \cdot (p \wedge u) = (c \cdot u)(d \cdot p) - (c \cdot p)(d \cdot u)$. This is $\mathcal{R}_{p}[h_{0}]$ up to the overall sign of the store's convention. Every statement below is about kernels, ranks, pairings or invariant subspaces, which an overall sign does not change. The Ricci map $\mathcal{R}_{p}(b) = \sum_{a} e^{a} \cdot \mathcal{R}_{p}(e_{a} \wedge b)$, the scalar, the Schouten map $\Sch(a) = \tfrac12(\mathcal{R}_{p}(a) - a\,\mathcal{R}_{p}/6)$ and $\mathcal{W}_{p}(a \wedge b) = \mathcal{R}_{p}(a \wedge b) - \Sch(a) \wedge b - a \wedge \Sch(b)$ follow; the check confirms that $\mathcal{W}_{p}$ is trace-free, which fixes the relative normalisation. A displacement $x \mapsto x + \epsilon\,\xi(x)$ acts at first order by $h(a) \mapsto h(a) + (a \cdot \nabla)\xi + \nabla(a \cdot \xi)$ (clause (e) of the definition, linearised); for $\xi(x) = \tfrac12 (p \cdot x)^{2}\xi_{0}$ the second derivatives change by $h_{0} \mapsto h_{0} + (a \cdot p)\xi_{0} + (a \cdot \xi_{0})p$, the space $G_{p}$. A conformal change $\hb^{-1}(a) = \phi\,a$ with $\phi = 1 + \tfrac14\sigma(p \cdot x)^{2}$ gives $h_{0} = \sigma\,\id$.

\emph{Covariance and representatives.} For a Lorentz map $\Lambda$, $\mathcal{R}_{\Lambda p}[\Lambda h_{0}\Lambda^{-1}] = \Lambda \mathcal{R}_{p}[h_{0}](\Lambda^{-1}\,\cdot\,)$ on bivectors, $G_{\Lambda p} = \Lambda G_{p}\Lambda^{-1}$, and $L$ is invariant; replacing $p$ by $cp$ multiplies $\mathcal{R}_{p}$ by $c^{2}$ and leaves $G_{p}$ unchanged. Reflections $a \mapsto -n a n^{-1}$ act in the same way. Every nonzero $p$ is so related to one of $e_{0}$, $e_{3}$ or $e_{0} + e_{1}$. Hence kernel dimensions, the inclusions below and the dimension of the stabiliser commutant are determined by these three momenta. The check computes them there and, as controls, at three further rational momenta of each type.

\emph{(i)} $\ker \mathcal{R}_{p} \supseteq G_{p}$ because a pure displacement has zero curvature; directly, $h_{0} = p \otimes \xi + \xi \otimes p$ gives $h_{0}(B \cdot p) = ((B \cdot p) \cdot \xi)\,p$ since $(B \cdot p) \cdot p = 0$, so $p \wedge h_{0}(B \cdot p) = 0$. The check finds $\dim\ker\mathcal{R}_{p} = 4 = \dim G_{p}$ at the representatives, so equality holds (computed exactly, proved for every $p$ by covariance). For $\mathcal{W}_{p}$: $G_{p} \subseteq \ker$ because $\mathcal{R}_{p}$ vanishes there, and $\id \in \ker$ because $\mathcal{R}_{p}[\id](B) = -\tfrac12 p \wedge (B \cdot p)$ is of the form $\Sch(a) \wedge b + a \wedge \Sch(b)$ (it is conformally flat; the check verifies $\mathcal{W}_{p}[\id] = 0$). The computed kernel has dimension $5 = \dim(G_{p} \oplus L)$ at all three representatives, null included. The scalar $h_{0} \mapsto \mathcal{R}_{p}[h_{0}]$ is a nonzero linear form (kernel of dimension $9$), vanishing on $G_{p}$, and its value on $\id$ is $3\,Q(p)$; for $Q(p) \neq 0$ its kernel $K_{p}$ therefore meets $G_{p} \oplus L$ in $G_{p}$, and we write $C_{p}$ for its image in $\mathrm{Sym}/G_{p}$, of dimension $5$.

\emph{(ii)} Let $H_{p}$ be the stabiliser of $p$. It preserves $G_{p}$, $L$ and $K_{p}$ (the latter because the scalar curvature is invariant), so the images $\bar L$ and $C_{p}$ of $L$ and $K_{p}$ are invariant subspaces of $\mathrm{Sym}/G_{p}$ with $\bar L \oplus C_{p} = \mathrm{Sym}/G_{p}$ (dimensions $1 + 5 = 6$). The check builds the action of five elements of $H_{e_{0}}$ (three rotations given by rational rotors $U$ with $\Lambda(a) = U a \rev{U}/(U\rev{U})$, rotation angles with cosines $\tfrac35$, $\tfrac{5}{13}$, $-\tfrac45$, and two reflections) and of five elements of $H_{e_{3}}$ (a rotation, two boosts from the rotors $2 + e_{0}e_{1}$ and $3 + e_{0}e_{2}$, with $\cosh = \tfrac53$ and $\tfrac54$, and two reflections) on $\mathrm{Sym}/G_{p}$, verifies that $G_{p}$ is invariant, and solves the commutation equations exactly: the commutant of these elements has dimension $2$. The commutant of $H_{p}$ is contained in it and contains the two projections onto $\bar L$ and $C_{p}$, so it equals their span. Each summand is irreducible under $H_{p}$: the stabiliser is $O(3)$ (compact) for timelike $p$ and $O(1,2)$ (semisimple identity component, finitely many components) for spacelike $p$, so its finite-dimensional real representations are completely reducible; a proper nonzero invariant subspace of $C_{p}$ would then have an invariant complement in $C_{p}$, and the projection onto it would be a third element of the commutant of $H_{p}$, independent of the first two, while that commutant lies in the computed two-dimensional commutant of the five elements. Hence the invariant subspaces of the quotient are $0$, $\bar L$, $C_{p}$ and the whole, since $\bar L$ and $C_{p}$ are inequivalent (different dimensions) and an invariant subspace of a sum of two inequivalent irreducibles is a sum of some of them. Lifting, the subspaces of $\mathrm{Sym}$ that contain $G_{p}$ and are $H_{p}$-invariant are $G_{p}$, $G_{p} \oplus L$, $K_{p} = G_{p} \oplus C_{p}$ (any lift) and $\mathrm{Sym}$. \emph{Null momentum.} At $p = n = e_{0} + e_{1}$ the stabiliser contains the transverse rotation from the rotor $2 - e_{2}e_{3}$, the null rotations from the rotors $1 + n e_{2}$ and $1 + n e_{3}$ (each has $U\rev{U} = 1$ since $Q(n) = 0$ and $n e_{i} = -e_{i}n$, and fixes $n$), and the reflection of $e_{2}$; it is not reductive and complete reducibility fails. With $N_{p}$ spanned by $a \mapsto (a \cdot e_{2})e_{2} - (a \cdot e_{3})e_{3}$ and $a \mapsto (a \cdot e_{2})e_{3} + (a \cdot e_{3})e_{2}$ (the helicity-two directions), the check finds $G_{p} \oplus N_{p}$ and $G_{p} \oplus L \oplus N_{p}$ invariant under these elements, $N_{p} \not\subseteq \ker\mathcal{W}_{p}$, and $L \subseteq K_{p}$; with $G_{p}$, $G_{p} \oplus L$, $K_{p}$ and $\mathrm{Sym}$ these are at least six invariant classes containing $G_{p}$. So the count of four holds only for $Q(p) \neq 0$. By (i), $G_{p} \oplus L = \ker\mathcal{W}_{p}$ is the first-order tangent space of the Weyl-flat fields, which by Theorem \claimref{synthesis.twistor-flat-fields-give-nordstrom-not-einstein}(i) are the twistor-flat fields.

\emph{(iii)} $T_{0}(p) = 0$ is four linear conditions of rank $4$, so the conserved sources form a space of dimension $6$ (computed). For $h_{0} = p \otimes \xi + \xi \otimes p$, $\ip{h_{0}}{T_{0}} = 2\,\xi \cdot T_{0}(p) = 0$. The check finds the pairing matrix of $\mathrm{Sym}$ against the conserved sources of rank $6$, so the induced pairing of the six-dimensional $\mathrm{Sym}/G_{p}$ with the conserved sources is nondegenerate. Since $\ip{\id}{T_{0}} = \tr T_{0} = \tau$, the conserved sources annihilating $G_{p} \oplus L$ are those with $\tau = 0$, a space of dimension $5$ (computed: the annihilator has dimension $5$ and consists of trace-free maps). For $Q(p) \neq 0$ the annihilator of $K_{p}$ has dimension $6 - 5 = 1$; the map $T_{0}(a) = Q(p)\,a - (a \cdot p)\,p$ is conserved ($T_{0}(p) = 0$) and invariant under $H_{p}$, and the check finds the computed annihilator equal to its span. A field equation from varying within a class $U \supseteq G_{p}$, $\ip{\delta h_{0}}{X} = 0$ for $\delta h_{0} \in U$ with $X = E_{0} + \kappa T_{0}$ conserved (the gravitational side conserved by the Noether identity of displacements, the matter side by its own equations), constrains $X$ exactly modulo the annihilator of $U$: for $U = G_{p} \oplus L$ it is the single condition $\tr X = 0$, and the trace-free part of $T_{0}$ (five components) does not enter, which is Nordstr\"om's coupling (Fact \claimref{synthesis.imported-nordstrom-scalar-theory-einstein-fokker-form}, $\mathcal{R} = -24\pi G\tau$ with $\mathcal{W} = 0$); for $U = \mathrm{Sym}$ it is $X = 0$, all six.

\emph{(iv)} The map $T \mapsto (k \cdot T(k))_{k}$ over twelve rational null vectors $k$ (among them $e_{0} \pm e_{i}$ and Pythagorean quadruples) has kernel of dimension $1$, spanned by $\id$, computed exactly; since $k \cdot \id(k) = Q(k) = 0$ the kernel always contains $\id$, and twelve directions already reach the bound, so the null fluxes determine $T$ modulo $\R\,\id$ for all null directions. If $X = \lambda\,\id$ is conserved at $p \neq 0$ then $X(p) = \lambda p = 0$ forces $\lambda = 0$ (the check also finds the conserved sources and $\id$ independent at three momenta). So for conserved $X$ at $p \neq 0$, vanishing trace-free part implies $X = 0$; at $p = 0$ a constant $\lambda\,\id$ survives, the cosmological term (the unimodular mechanism).

\emph{(v)} Let $\mathcal{L}(\hb) = (\det\hb)^{-1}\Phi(\hb(\partial\psi), \psi)$ where $\hb$ enters only through the derivatives of the matter field (applied as a linear map, or as its outermorphism on a field strength) and the volume factor. Along $\hb = \id + t\,\delta$: $\frac{d}{dt}(\det\hb)^{-1}|_{0} = -\tr\delta$ and $\frac{d}{dt}\Phi|_{0}$ is, by the chain rule, $\partial\Phi/\partial(\nabla\psi)$ contracted with $\delta(\nabla\psi)$, so the derivative is linear in $\delta$, of the form $\sum_{a} \delta(e_{a}) \cdot T(e^{a})$, with $T(a)$ the contraction of $\partial\Phi/\partial(\nabla\psi)$ with $a \cdot \nabla\psi$ minus $a\,\mathcal{L}$, the $\hb$-response stress. When $\Phi$ depends on $\nabla\psi$ directly (the scalar) this is the canonical Noether stress of displacements (the current of global translations); when $\hb$ acts on a field strength built from a potential (the Maxwell field, $F = \nabla \wedge A$ held fixed) it is the symmetric gauge-invariant stress, which differs from the canonical Noether stress of $A$ by the Belinfante term. For the scalar, $\Phi = \tfrac12 Q(\hb(\nabla\phi)) - \tfrac12 m^{2}\phi^{2}$ gives $T(a) = (a \cdot \nabla\phi)\nabla\phi - a\,\mathcal{L}$ and $\tr T = Q(\nabla\phi) - 4\mathcal{L} = -Q(\nabla\phi) + 2m^{2}\phi^{2}$. For the Maxwell field, $\Phi = \tfrac12\br{\hb(F)^{2}}$ with $\hb(a \wedge b) = \hb(a) \wedge \hb(b)$; at first order $\hb(F) = F + t\,\delta_{\wedge}F$ with $\delta_{\wedge}(a \wedge b) = \delta(a) \wedge b + a \wedge \delta(b)$, so the derivative is $\br{F\,\delta_{\wedge}F} - \tfrac12\tr\delta\,\br{F^{2}}$, which equals $\sum_{a}\delta(e_{a}) \cdot T(e^{a})$ with $T(a) = -\tfrac12 FaF$, the stress of Fact \claimref{synthesis.imported-maxwell-stress-spacetime-algebra}, symmetric and gauge invariant (the Belinfante-improved canonical stress). The check does not use these formulas: it evaluates the full nonlinear densities exactly at six rational $t$ for each of the sixteen elementary $\delta$ and takes the derivative at $0$ of the interpolating polynomials (both numerator and $\det\hb$ are polynomials of degree at most $4$ in $t$), then compares with the two stress maps at a rational field strength and a rational gradient; it also confirms that the Maxwell response is symmetric and trace-free and the scalar trace formula. For $\delta = \id$, $\hb = (1 + t)\id$ multiplies $\hb(F)$ by $(1 + t)^{2}$ and $\det\hb$ by $(1 + t)^{4}$, so the Maxwell density is constant in $t$: a density algebraic in $\hb$ and invariant under constant rescaling has zero response along $\id$, at each point. A first-order Weyl-flat variation is $\delta = \sigma(x)\,\id$ plus a displacement; the displacement part pairs to zero with the conserved stress after integration by parts, and the conformal part gives $\int \sigma\,\tr T = 0$. So the Maxwell field does not couple at first order to variations within the Weyl-flat class. For the Dirac field the response is not algebraic in $\hb$ when $\Omega$ is eliminated, and with $\Omega$ independent it is the canonical stress (see (vii)); the classical on-shell trace of its stress is proportional to $m$ times the scalar spin product, so for $m = 0$ it vanishes on shell (not pointwise off shell, as for the Maxwell density), by conformal invariance with weight $-\tfrac32$ for the spinor (standard, cited from memory, not checked). At the quantum level the renormalised stress has a state-independent trace anomaly containing a scheme-dependent $\Box\mathcal{R}$ term, which is first order in $h$ and nonzero on $L$ (the scalar of $\id$ is $3\,Q(p)$), so the decoupling of the massless Dirac field from the Weyl-flat class is classical only.

\emph{(vi)} For point sources with momenta $p_{n}$ coupled to the field at a soft momentum $q$ through $\kappa_{n}\,p_{n} \cdot h_{0}(p_{n})/(p_{n} \cdot q)$ (the classical radiated field of a scattering at a point, or Weinberg's soft factor), the change under $h_{0} \mapsto h_{0} + q \otimes \xi + \xi \otimes q$ is $2\,\xi \cdot \sum_{n} s_{n}\kappa_{n}p_{n}$, which must vanish for every $\xi$. On the configuration of the check (five rational momenta, $\sum s_{n}p_{n} = 0$, spanning $V$, so one linear relation among them) the solution space of $\sum s_{n}\kappa_{n}p_{n} = 0$ is computed to be the line of constant $\kappa$: one coupling for all five sources. With six or more momenta a single configuration leaves a kernel of dimension $n - 4$, and universality follows only by imposing the condition on every momentum-conserving configuration, which is Weinberg's argument; that general step is standard and not computed here. Within $G_{p} \oplus L$ the coupling of a source is $\sigma\,p_{n} \cdot p_{n} = \sigma\,Q(p_{n})$, zero for a massless source (the stored absence of light deflection in Nordstr\"om's theory).

\emph{(vii), argued.} By (i) and (ii) the Weyl-flat class is one of the four displacement-closed covariant classes, so displacement gauge invariance alone does not exclude it, and by (iii) it couples to the trace only; the Einstein-Fokker form of Nordstr\"om's theory is generally covariant, the nonlinear counterpart. If the gravitational field equations are the Euler-Lagrange equations of an action $S_{g} + S_{m}$ with $S_{m}$ minimally coupled and $\hb$ varied in all directions, the source of the $\hb$ equation is the response of (v), and by (iii) and (iv) varying in the directions $C_{p}$ supplies its trace-free part, which with conservation fixes it at every nonzero momentum. With Lovelock's theorem (Fact \claimref{synthesis.imported-lovelock-and-anderson-duchamp-natural-lagrangians}(a)) a natural divergence-free second-order gravitational side is $c_{1}\mathcal{G} + c_{2}$, which gives Einstein's equations. The directions $C_{p}$ leave $\ker\mathcal{W}_{p}$, the tangent of the attachment's domain, so the clause holds under that action principle exactly when the matter theory is defined and varied there. For matter with spin, the definition's note makes $\Omega$ independent; then the $\hb$-response of (v) is the canonical stress, whose antisymmetric part equals the divergence of the spin density on shell; it enters the antisymmetric part of the $\hb$ equation, and the $\Omega$ equation ties the torsion to the spin density (Einstein-Cartan, standard, cited from memory). Eliminating the torsion, the symmetric source is the Belinfante stress plus a spin-spin contact term of order $\kappa$, which is second order in the coupling; at first order the pairings of (iii), which involve only symmetric $h_{0}$, therefore see the Belinfante stress. Conservation of the source, used in (iii) and (iv), must survive quantisation; for chiral matter that is the cancellation of the mixed gravitational anomaly (argued, not used). This step is a reading of (i) to (vi) at first order; the nonlinear statement and the action principle itself are not derived.

\emph{Circularity.} The uses are the twistor-flat result (for the identification of the Weyl-flat class with the attachment's domain), the Maxwell stress Fact, Nordstr\"om's Fact, Lovelock's Fact and the dictionary Fact; none assumes that the whole stress is the source. Einstein's equations enter only in (vii), through Lovelock's theorem, after the clause.
